\documentclass[11pt]{article}
\usepackage[margin=1in]{geometry}
\usepackage{booktabs}
\title{A Repaired Exact Laboratory for Adversarial Boolean-Network Repertoires}
\author{CausalBool group}
\date{2026-09-09}

\begin{document}
\maketitle

\begin{abstract}
We define and validate a finite Boolean-network laboratory for studying how
small wiring perturbations constrain attainable long-run distributions and
payoffs. The repaired contract uses row-input adjacency, an explicit
basin-weighted attractor repertoire, standard forward KL, and a canonical
reversible distribution codec. An exact N=6 study covers four topology
families, two independent base seeds, and separate edge additions and
removals. The resulting artefacts are reproducible and validated; this report
does not claim large-network scaling or detector performance.
\end{abstract}

\section{Contract}
The adjacency convention is $A[\mathrm{target}][\mathrm{source}]=1$. We record
the one-step transition map, its uniform-input output distribution, and the
primary long-run distribution. The latter follows every uniformly sampled
initial state to its deterministic cycle and phase-averages that cycle. For
distributions on the full $N$-bit alphabet, $D_{KL}(p\Vert q)$ is standard
forward KL and is infinite when $p$ has mass outside the support of $q$.
Conditional KL is secondary. Additions and removals are separate estimands .

\section{Exact study}
The prespecified run uses ring, sparse-random, modular, and hub families; two
independent seeds; maximum in-degree three; no empty inputs; no self-loops;
$k=1$; a single baseline-support target selected by seed; and $C=0.25$.
There are 16 family/seed/operation runs and 219 admissible perturbation rows.
All 219 rows pass the independent trajectory oracle and have normal process
status. Ninety-one rows have infinite forward KL, demonstrating why support
intersection must not replace the declared alarm statistic. Exclusions and
failures are retained in each run summary.

The finite-catalogue value is computed as
\[
 V_k(C)=\max_{A:\,D_{KL}(\mu_A\Vert q)\le C} E_{\mu_A}[s].
\]
The exponential-family KL-ball expression is reported only as an upper-bound
benchmark; its gap from $V_k(C)$ measures the restriction imposed by the
network catalogue. It is not asserted that the exponential tilt is a network
attacker.

\section{Reproducibility and limitations}
Run the experiment from the repository root with
\begin{verbatim}
PYTHONPATH=doppel-challenge/src python -m doppel_challenge \\
  --out-dir doppel-challenge/results/exact_small
\end{verbatim}
Each case has an atomic checkpoint, a stable perturbation identifier, source
and environment provenance, and a scientific digest. The descriptive report
is generated only from accepted rows. The historical pilot files are retained
as exploratory diagnostics and are not used in the denominators above.

The experiment is finite-state and exact only at the declared small size. It
does not identify a wiring mechanism from a distribution when distinct
networks collide, and it does not establish a finite-observation detector.
Held-out benign/attack sampling, false-positive calibration, uncertainty
under clustered perturbations, and approximate larger-$N$ estimators require a
separate declared study.

\section*{Data and software}
The protocol, schema, code, and hashed artefacts are in
\texttt{doppel-challenge/}. The committed exact-study summary is
\texttt{results/exact\_small/study\_summary.json}; its descriptive analysis is
\texttt{results/exact\_small/analysis.json}.

\end{document}
